\fsection{Necesidades de automatización identificadas}
Del análisis del entorno se identifican las siguientes necesidades principales:

\begin{itemize}
	\item Clasificación y trazabilidad de producto: empresas logísticas y agroalimentarias necesitan identificar y separar productos por tipo, lote o destino de forma automatizada, reduciendo errores y costes de mano de obra.

	\item Picking flexible: la variabilidad de los productos manipulados (distintas formas, tamaños, pesos) hace inviables los sistemas rígidos de posicionamiento previo. Los sistemas de bin picking con visión artificial son una solución directa a esta necesidad.

	\item Integración de trazabilidad RFID: la normativa de trazabilidad en el sector alimentario y la creciente exigencia de los operadores logísticos impulsan la adopción de sistemas de identificación automática.
\end{itemize}
